\documentclass[11pt]{article}

\usepackage[a4paper,margin=2.4cm]{geometry}
\usepackage{amsmath,amssymb}
\usepackage{booktabs}
\usepackage{listings}
\usepackage{xcolor}
\usepackage[colorlinks=true,linkcolor=blue,urlcolor=blue]{hyperref}

\lstset{
  basicstyle=\ttfamily\footnotesize,
  breaklines=true,
  columns=fullflexible,
  keepspaces=true,
  frame=single,
  framesep=4pt,
  xleftmargin=4pt,
}

\newcommand{\g}{g}
\newcommand{\Tr}{\operatorname{Tr}}
\newcommand{\bt}{\beta}

\title{Exploiting the imaginary-time reflection symmetry in spinDMFT}
\author{spinDMFT (Matsubara / pCN) algorithm note}
\date{\today}

\begin{document}
\maketitle

\section{The symmetry}

The imaginary-time spin correlation
\begin{equation}
  \g^{ab}(\tau) \;=\; \big\langle\, S^a(\tau)\, S^b(0)\,\big\rangle
\end{equation}
obeys, for the equilibrium ensemble,
\begin{equation}
  \boxed{\;\g^{ab}(\bt-\tau) \;=\; \g^{ab}(\tau)^{*}\;}
  \qquad\Longleftrightarrow\qquad
  \begin{aligned}
    \operatorname{Re}\g^{ab}(\bt-\tau) &= \phantom{-}\operatorname{Re}\g^{ab}(\tau),\\
    \operatorname{Im}\g^{ab}(\bt-\tau) &= -\operatorname{Im}\g^{ab}(\tau).
  \end{aligned}
  \label{eq:sym}
\end{equation}
This is a \emph{same-component} relation: each $\g^{ab}$ has an even real part and an
\emph{odd} imaginary part about $\bt/2$. It follows from two exact identities,
\begin{equation}
  \g^{ab}(\bt-\tau) = \g^{ba}(\tau)\quad\text{(trace cyclicity)},
  \qquad
  \g^{ab}(\tau)^{*} = \g^{ba}(\tau)\quad\text{(Hermiticity)},
  \label{eq:twoid}
\end{equation}
whose composition gives \eqref{eq:sym}. \textbf{Caveat:} combining them as
$\g^{ab}(\bt-\tau)=\g^{ba}(\tau)^{*}$ is wrong --- it double-applies the conjugate
($\g^{ba}(\tau)^{*}=\g^{ab}(\tau)$) and collapses to
$\g^{ab}(\bt-\tau)=\g^{ab}(\tau)$, which makes the imaginary part \emph{even}. The
error is invisible on diagonal ($a=b$) components, where the transpose is a no-op, and
appears only for off-diagonal components (e.g.\ in an external field).

\section{Origin: it is an \emph{ensemble} property}

In spinDMFT each Monte-Carlo sample is a noise trajectory $V(\tau)$, and the
single-sample correlation is
\begin{equation}
  \g^{ab}(\tau;V) \;=\; \frac{1}{Z(V)}\,
  \Tr\!\big[\,U(\bt,\tau)\,S^a\,U(\tau,0)\,S^b\,\big],
  \qquad
  U(\tau_2,\tau_1)=\mathcal{T}\,e^{-\int_{\tau_1}^{\tau_2} H(s)\,ds},
\end{equation}
with $Z(V)=\Tr U(\bt,0)$. Because $U$ is a time-ordered product of exponentials of
\emph{non-commuting} Hermitian generators, it is \emph{not} itself Hermitian, and the
cyclicity identity in \eqref{eq:twoid} does \emph{not} hold sample by sample: writing
$U(\bt,0)=U_{N-1}\cdots U_1U_0$ with $U_k=e^{-\delta H(\tau_k)}$, the order of the
factors flanking $S^a$ and $S^b$ cannot be reversed for a generic trajectory. It is
restored only after averaging, because the noise distribution is invariant under
reflection $V_R(\tau)\equiv V(\bt-\tau)$, $p(V)=p(V_R)$ (the covariance is even / KMS,
the mean field is constant):
\begin{equation}
  \langle \g^{ab}(\bt-\tau)\rangle
  = \int\! p(V)\,\g^{ab}(\bt-\tau;V)
  = \int\! p(V_R)\,\g^{ba}(\tau;V_R)
  = \langle \g^{ba}(\tau)\rangle .
\end{equation}
Together with the (likewise ensemble-level) Hermiticity identity this yields
\eqref{eq:sym}.

\paragraph{Consequence.} A strategy that \emph{enforces $V(\bt-\tau)=V(\tau)$ per
sample} restricts the sampler to the symmetric (even-frequency) subspace. That is a
\emph{different, biased} distribution, not $p(V)$. It was tried and rejected: it
agrees with the reference at $\bt=1$ but deviates by $7$--$9\,\sigma$ at $\bt=2,3$.
The correct way to use \eqref{eq:sym} is therefore on the \emph{ensemble average}, not
on the field.

\section{Implemented strategy: half-range computation + reconstruction}

On the grid $\tau_k=k\,\delta t$, $k=0,\dots,N_{\mathrm{ts}}$ with
$\delta t=\bt/N_{\mathrm{ts}}$ and $N=N_{\mathrm{ts}}+1$ time points, the reflection
$\bt-\tau_k=\tau_{N-1-k}$ maps grid point $k$ to point $N-1-k$. We therefore

\begin{enumerate}
  \item sample the per-sample correlation only on the lower half
        $k=0,\dots,n_{\mathrm{half}}-1$ with $n_{\mathrm{half}}=N-\lfloor N/2\rfloor$
        (the points in $[0,\bt/2]$), and
  \item once per self-consistent iteration, after the Monte-Carlo average and the
        standard-error estimate, reconstruct the upper half from \eqref{eq:sym}:
        \[
          \operatorname{Re}\g^{ab}_k =  \operatorname{Re}\g^{ab}_{N-1-k},\quad
          \operatorname{Im}\g^{ab}_k = -\operatorname{Im}\g^{ab}_{N-1-k},
          \qquad k\ge n_{\mathrm{half}} .
        \]
\end{enumerate}
The mirror partner $N-1-k$ always lies in the computed lower half. The mapping is
\emph{per component} --- no transpose --- so each $\g^{ab}$ is reconstructed from itself
and no cross-pair bookkeeping is needed. The per-point standard errors are reconstructed
by the same mirror with \emph{no} sign change, because each reconstructed value equals
its mirror source deterministically. The self-paired midpoint $\tau=\bt/2$ (present only
for an odd number of time points) lies in the computed lower half and keeps its raw
sampled value: its imaginary part vanishes only on the ensemble average, so forcing it
to zero per iteration would bias the estimator rather than improve it.

This is unbiased (the lower half is an ordinary MC estimate, the upper half follows
exactly), enforces the symmetry to machine precision, and roughly halves the dominant
per-sample observable work --- the $S^a(\tau)$ operator products and the correlation
traces. The full propagator chain is still built over $[0,\bt]$ because the partition
function $Z=\Tr U(\bt,0)$ needs it, so the wall-clock speedup is partial rather than
exactly $2\times$.

\section{Code changes}

\begin{description}
  \item[\texttt{Functions.cpp} -- \texttt{compute\_S\_of\_t}] builds the $S^a(\tau)$
        caches only for the $n_{\mathrm{half}}$ lower-half points (the propagator loop
        is truncated at \texttt{propagators.cbegin() + n\_half}).
  \item[\texttt{Functions.cpp} -- \texttt{compute\_spin\_correlations}] the inner
        accumulation loop is now driven by the (half-sized) $S^a(\tau)$ caches, so the
        sums and squared-sums are filled only on $[0,\bt/2]$.
  \item[\texttt{Functions.cpp} -- \texttt{symmetrize\_upper\_half} (new)]
        reconstructs the upper-half correlations and their standard errors from the
        lower half using \eqref{eq:sym}.
  \item[\texttt{main.cpp}] calls \texttt{symmetrize\_upper\_half} after
        \texttt{compute\_sample\_stds} and before the convergence check, so the
        fed-back correlator and the iteration-error metric both see the full curve.
\end{description}

\begin{lstlisting}[caption={Reconstruction kernel (Functions.cpp).}]
void symmetrize_upper_half( CorrTen& Re, CorrTen& Im,
                            CorrTen& Re_std, CorrTen& Im_std )
{
    const size_t n_pairs = Re.size();
    const size_t N = Re[0].size();
    const size_t n_half = N - N / 2;            // computed points in [0, beta/2]
    for( size_t p = 0; p < n_pairs; ++p )
    {
        for( size_t k = n_half; k < N; ++k )
        {
            const size_t m = N - 1 - k;          // mirror index (lower half)
            Re[p][k]     =  Re[p][m];            // same component: Re even
            Im[p][k]     = -Im[p][m];            //                 Im odd
            Re_std[p][k] =  Re_std[p][m];
            Im_std[p][k] =  Im_std[p][m];
        }
        // midpoint tau = beta/2 (odd N) keeps its raw sampled value
    }
}
\end{lstlisting}

\section{Verification}

Diagonal correlator $\g^{aa}_{ii}(\tau)$, ISO model, spin-$1/2$, $J_Q=1$,
$N_{\mathrm{ts}}=49$, seed fixed; half-range runs use $4\times10^4$ samples. The
references are the largest available standard runs. ``max/std'' is the largest
pointwise deviation from the reference in units of the half-range standard error.

\begin{center}
\begin{tabular}{cccccc}
\toprule
$\bt$ & reference & symmetry residual & $\max|\text{half}-\text{ref}|$ & half std (mean) & max/std \\
\midrule
1 & $4\times10^{6}$ & $0.0$ & $7.3\times10^{-5}$ & $8.9\times10^{-5}$ & $0.82$ \\
2 & $8\times10^{4}$ & $0.0$ & $3.3\times10^{-4}$ & $7.8\times10^{-4}$ & $0.43$ \\
3 & $4\times10^{6}$ & $0.0$ & $1.4\times10^{-3}$ & $1.2\times10^{-3}$ & $1.16$ \\
\bottomrule
\end{tabular}
\end{center}

The result agrees with the reference within one standard error at every $\bt$
(max/std $\approx 1$), and the symmetry residual
$\max_\tau|\operatorname{Re}\g(\tau)-\operatorname{Re}\g(\bt-\tau)|$ is exactly zero
because the upper half is defined by the reconstruction rather than sampled.

\paragraph{Off-diagonal check (external field).} The ISO benchmark above only exercises
the diagonal components, where the (incorrect) transposed reconstruction happens to give
the right answer. To exercise the off-diagonals we run $\bt=3$ with a field $h_z=0.5$ and
symmetry type C (rows $[xx,xy,yx,zz]$), which turns on $\g^{xy},\g^{yx}$ with large
imaginary parts. With \eqref{eq:sym} the per-component residuals
$\max_\tau|\operatorname{Re}\g^{ab}(\tau)-\operatorname{Re}\g^{ab}(\bt-\tau)|$ and
$\max_\tau|\operatorname{Im}\g^{ab}(\tau)+\operatorname{Im}\g^{ab}(\bt-\tau)|$ are all
zero, and $\operatorname{Im}\g^{xy}$ runs $-0.13 \to 0 \to +0.13$ across $[0,\bt]$ ---
i.e.\ genuinely \emph{odd}. The earlier (transposed) kernel instead produced an even
$\operatorname{Im}\g^{xy}$, which is unphysical.

\end{document}
